% answers_short.tex —— 2008 年·普通高等学校招生考试·福建卷·理科 · 简短答案速查
% 由 tools/make_short_answers.py 自 answers.tex 的 \gksolhead/\gkanswer 提取，
% 并与 2008/2008_数学真题_福建理_答案要点.md 逐题比对后生成，勿手改（改源 MD 后重跑）。

\begin{gknotice}
\textbf{使用说明}：本节只给最终答案便于快速核对，完整推导见解析版。答案为编写者独立推导，\textbf{非官方标准答案}。
\end{gknotice}

\gkpart{选择题}
\begin{center}\begin{tabular}{ccccccccccccc}
\toprule
题号 & 1 & 2 & 3 & 4 & 5 & 6 & 7 & 8 & 9 & 10 & 11 & 12 \\
\midrule
答案 & B & A & C & B & B & D & A & C & A & D & B & C \\
\bottomrule
\end{tabular}\end{center}

\gkpart{填空题}
\begin{itemize}[leftmargin=2.6em]
  \item[\textbf{13.}] $31$。
  \item[\textbf{14.}] $(-\infty,0)\cup(10,+\infty)$。
  \item[\textbf{15.}] $9\pi$。
  \item[\textbf{16.}] ③④。
\end{itemize}

\gkpart{解答题}
\begin{itemize}[leftmargin=2.6em]
  \item[\textbf{17.}] （1）$A=\dfrac{\pi}3$；（2）值域 $\left[-3,\dfrac32\right]$。
  \item[\textbf{18.}] （1）证明见解析；（2）$60^\circ$；（3）存在，$\dfrac{AQ}{QD}=\dfrac13$。
  \item[\textbf{19.}] （1）证明见解析；（2）$-2<a<-1$ 时极大值 $-\dfrac23$；$0<a<1$ 时极小值 $-2$；其余情形无极值。
  \item[\textbf{20.}] （1）$\dfrac13$；（2）$E\xi=\dfrac83$。
  \item[\textbf{21.}] （1）$\dfrac{x^2}4+\dfrac{y^2}3=1$；（2）$a>\dfrac{1+\sqrt5}2$。
  \item[\textbf{22.}] （1）增区间 $(-1,0)$，减区间 $(0,+\infty)$；（2）① $c\leqslant1$；② 证明见解析。
\end{itemize}
